\documentclass{article}
\usepackage{iclr2027_conference,times}
\usepackage{amsmath,amssymb}
\usepackage{booktabs}
\usepackage{enumitem}
\usepackage{graphicx}
\usepackage{hyperref}
\usepackage{url}

% \iclrfinalcopy  % enable for camera-ready

\title{Task-efficient agent evaluation using cached traces}

\author{Anonymous}

\newcommand{\trubric}{\textsc{trubric}}

% Notation: systems f in F, benchmark Q* = {q_1,...,q_M}, trace-level scorer s,
% benchmark score y, references (f_i, y_i), budget m << M. Artifacts live in
% ../artifacts/ (trubric_table1_hero.tex, trubric_figure{1..6}_*.pdf).

\begin{document}
\maketitle

\begin{abstract}
Efficient benchmarking estimates a system's full-benchmark performance from a small number of task executions.
Standard approaches rely solely on the binary correctness of executions from a collection of already-evaluated reference systems.
An agentic system, however, produces a rich trace for each execution that can reveal substantially more about its capabilities.
Exploiting this information requires useful trace representations.
We propose representing a trace by extracting a structured, rubric-guided summary of it with an LLM and embedding the summary with a standard text encoder.
These representations preserve correctness- and task-related signal while minimizing the presence of system identity-, model family-, and harness-related signals that dominate raw trace embeddings.
Given this representation for each task execution from all reference systems, we construct low-dimensional representations of the systems, predict a new system's score from its position among the reference systems, and combine the resulting geometric prediction with a two-parameter logistic item-response model.
On SWE-bench Verified and Terminal-Bench 2.0, the proposed estimator reduces mean absolute error from 0.13 to 0.08 ($p < 10^{-4}$) and from 0.14 to 0.12 ($p < 10^{-3}$), respectively, when provided a single random executed task relative to an item-response model alone.
These results indicate that trace geometry and correctness outcomes carry complementary predictive signal and enable more accurate benchmark score estimation from limited agent executions.
\end{abstract}

\section{Introduction}
\label{sec:intro}

The dominant paradigm of generative model use is increasingly agentic: systems plan, invoke tools, observe intermediate results, and iterate over long horizons to complete complex tasks \citep{yang2024sweagent, wang2025openhands}.
As agentic capabilities increase, so does the range of settings in which a given system might be deployed -- and, with it, the number of benchmarks on which the system must be evaluated before it can be trusted with consequential work \citep{kapoor2025hal}.
Agentic evaluation, however, is expensive as each task execution requires a properly sandboxed environment and a single task execution routinely consumes hundreds of thousands of tokens \citep{jimenez2024swebench, tbench2026} and can regularly exceed \$2{,}000 \citep{kapoor2025hal}.
% At reported per-task costs of \$0.13--\$4 \citep{kapoor2025hal}, the cost of evaluating a single system on a single benchmark can regularly exceed \$2{,}000.
Given this cost, efficient evaluation is necessary to increase the evaluation coverage of an agentic system.

Task-efficient evaluation methods (e.g., benchmark compression and anchor-item selection \citep{vivek2024anchor, maiapolo2024tinybenchmarks}, item response theory and adaptive testing \citep{kipnis2025metabench, truong2025amortized}, leveraging cached responses \citep{helm2026query}, etc.) address this problem by finding structure in item-level information available from previously-evaluated reference systems to decrease the cost of evaluating a new system.
Some task-efficient evaluation methods require only the item-level scores from reference systems \citep{vivek2024anchor, kipnis2025metabench}.
Others leverage the content of cached artifacts, such as natural language responses or residual streams, to measure the similarity of systems on a set of tasks \citep{duderstadt2023data, helm2026jailbreak, helm2026query}.
Notably, agentic workflows often produce artifacts when executing a task -- file diffs, tool calls, intermediate plans, and error messages -- recorded together in an execution trace that can be used to this effect.
Unfortunately, existing methods for using cached artifacts rely on off-the-shelf embedding functions that are not immediately well-suited for representing traces.

We make two primary contributions:
\begin{enumerate}
[wide, labelwidth=1pt, labelindent=12pt]
\setlength{\itemsep}{1pt}\setlength{\parskip}{0pt}
    \item \textbf{\trubric{} trace representations.} We introduce \trubric{} (\emph{trace rubric}): a trace representation constructed by extracting a task specific structured summary of each trace with an LLM and embedding the summary with a standard text encoder.
    We show that \trubric{} embeddings preserve information relevant to the task and behavior of the system while minimizing system-, model family-, and harness- specific signal across seven modern text encoders.
    \item \textbf{Task-efficient evaluation using cached traces.} Given the \trubric{} embeddings from a set of tasks for a target system and a collection of reference systems, we evaluate their utility for predicting the target system's full-benchmark score on SWE-bench Verified and Terminal-Bench 2.0 under three different agentic system evaluation protocols.
    We find that \trubric{}-based estimators, combined with an item-response model, reduce single-task mean absolute error from 0.13 to 0.08 ($p < 10^{-4}$) and from 0.14 to 0.12 ($p < 10^{-3}$) on the two benchmarks, respectively, and are nearly always better than baselines at any task budget (Table \ref{tab:hero}).
\end{enumerate}

% requires booktabs + graphicx (and the \trubric macro)
\begin{table}[t]
\centering
\caption{Mean Absolute Error (MAE) of benchmark score estimators for SWE-bench Verified and Terminal-Bench 2.0 under leave-one-family-out evaluation and different task selection protocols (random \& adaptive). \textbf{Bolded} values are the best method for a given benchmark, task selection, and number of tasks. \underline{Underlined} values are statistically significant ($\alpha = 0.05$) comparisons to IRT per the two-sided Wilcoxon signed-rank test. The \trubric{} and IRT ensemble -- which takes advantage of both trace geometry and per-execution score -- is the best or near-best in all settings.}
\label{tab:hero}
\resizebox{\textwidth}{!}{%
\begin{tabular}{lcccccccccccc}
\toprule
 & \multicolumn{6}{c}{\textbf{SWE-bench Verified}} & \multicolumn{6}{c}{\textbf{Terminal-Bench 2.0}} \\
\cmidrule(lr){2-7} \cmidrule(lr){8-13}
 & \multicolumn{3}{c}{random} & \multicolumn{3}{c}{adaptive} & \multicolumn{3}{c}{random} & \multicolumn{3}{c}{adaptive} \\
\cmidrule(lr){2-4} \cmidrule(lr){5-7} \cmidrule(lr){8-10} \cmidrule(lr){11-13}
\textit{Num.\ tasks $m$} & {1} & {5} & {20} & {1} & {5} & {20} & {1} & {5} & {20} & {1} & {5} & {20} \\
\midrule
sample score & 0.442 & 0.175 & 0.082 & 0.380 & 0.151 & 0.072 & 0.396 & 0.164 & 0.073 & 0.360 & 0.124 & 0.132 \\
IRT & 0.127 & 0.088 & 0.048 & 0.080 & 0.062 & \bfseries 0.034 & 0.137 & 0.097 & 0.052 & 0.112 & 0.067 & 0.039 \\
raw trace & {\underline{0.099}} & 0.093 & 0.089 & 0.096 & 0.083 & 0.081 & 0.135 & 0.126 & 0.120 & 0.128 & 0.125 & 0.113 \\
generic & {\underline{\bfseries 0.079}} & {\underline{0.071}} & 0.068 & 0.074 & 0.070 & 0.069 & 0.139 & 0.127 & 0.115 & 0.144 & 0.130 & 0.117 \\
trubric & {\underline{0.081}} & {\underline{0.069}} & 0.060 & 0.068 & 0.056 & 0.060 & {\underline{0.127}} & 0.107 & 0.093 & 0.110 & 0.092 & 0.084 \\
IRT $+$ trubric & {\underline{0.081}} & {\underline{\bfseries 0.063}} & {\underline{\bfseries 0.043}} & {\underline{\bfseries 0.062}} & {\underline{\bfseries 0.048}} & \bfseries 0.034 & {\underline{\bfseries 0.122}} & {\underline{\bfseries 0.090}} & {\underline{\bfseries 0.050}} & {\underline{\bfseries 0.099}} & {\underline{\bfseries 0.059}} & \bfseries 0.037 \\
\bottomrule
\end{tabular}}
\end{table}


\subsection{Problem statement}
\label{sec:problem}

For our purposes, an agentic system $f \in \mathcal{F}$ is a random mapping from a task space $\mathcal{Q}$ to a trace space $\mathcal{X}$.
A benchmark is a fixed collection of tasks $Q^{*} = \{q_{1}, \hdots, q_{M}\} \subset \mathcal{Q}$.
Each trace is evaluated by the benchmark's trace-level scoring function $s: \mathcal{X} \to \{0, 1\}$ and the system is evaluated by its benchmark score $ y(f) = \frac{1}{M} \sum_{j=1}^{M} \mathbb{E}\left[ s(f(q_{j})) \right] $
where the expectation is with respect to the randomness of $f$.
We write $y$ for $y(f)$ when the context is clear.

When evaluating a new system, we typically have access to $ \mathcal{D} = \left\{ \{ (x_{ij}, y_{ij}) \}_{j=1}^{M} \right\}_{i=1}^{n} $, the traces and trace-level scores of $n$ previously-evaluated reference systems $f_{1}, \hdots, f_{n}$,
where $x_{ij} \sim f_{i}(q_{j})$, $y_{ij} = s(x_{ij})$ and $y_{i} = y(f_{i}) $.
For the target system $f_{0}$, we observe traces (and, optionally, trace-level scores) on only a subset $Q \subset Q^{*}$ of $|Q| = m \ll M$ tasks and our goal is to estimate the system's full benchmark score $ y_{0} $ using only the target's $m$ executions and $\mathcal{D}$.

\section{Method}
\label{sec:method}
% TODO: trubric construction (rubric-guided extraction, embedding, consensus
% centering, PKPS geometry, estimators, IRT blend).

\section{Trace representations encode confounders}
\label{sec:probes}
% TODO: linear-probe analysis (Figure 1), PKPS geometry (Figure 2).

\section{Efficient evaluation}
\label{sec:results}
% TODO: protocol grid (Figure 3), ablations (Figure 4),
% cost and ranking (Figure 5), rubric sensitivity (Figure 6).

\section{Related work}
\label{sec:related}

\paragraph{Representing traces and structured artifacts.}
The standard tools for embedding model outputs are text encoders trained on sentence-level semantic similarity \citep{reimers2019sbert, su2023instructor, nussbaum2025trainingsparsemixtureexperts}.
Agentic traces, however, implicitly or explicitly encode the scaffold or underlying model that may not be relevant for downstream tasks \citep{sun2025idiosyncrasies, oderinwale2026trajectories} and are not necessarily suitable to be encoded.
Instead, traces have primarily been studied through taxonomies of failure modes and actions \citep{cemri2025mast, shu2026traceprobe}, LLM judges of agent behavior \citep{zhuge2025agentjudge, mehtiyev2026beyond}, and trajectory compression for cost reduction \citep{xiao2025agentdiet}.
Further, parallel work adapts model outputs to better fit the notion of semantic similarity either by comparing structured texts through natural-language descriptions of their properties \citep{ravfogel2024description} or by representing an output by an LLM's answers to a set of questions about it \citep{peng2024inbedder, benara2024qaemb}.
\trubric{} applies this to agentic traces by embedding a task-specific structured summary of each trace, preserving behavior-relevant information while minimizing system-level information.

\paragraph{Efficient model evaluation.}
Task-efficient evaluation methods estimate a system's full-benchmark score from a small number of executions.
Item response theory (IRT) is the workhorse of score-only approaches: each system is assigned a latent ability, each task a difficulty and discrimination, and a logistic link models the probability of success, so that an ability estimated from $m$ graded executions yields a predicted benchmark score \citep{maiapolo2024tinybenchmarks, kipnis2025metabench}.
Extensions select informative or adaptive task subsets \citep{vivek2024anchor, truong2025amortized} and apply IRT to agentic benchmarks directly \citep{ge2026psychometrics, ndzomga2026efficient}, though the reliability of small-budget estimates remains contested \citep{yauney2026reliable}.
A complementary family represents the systems themselves, either from their score profiles \citep{ruan2024observational, zhuang2025embedllm} or from the content of their cached responses.
The Data Kernel Perspective Space (DKPS) embeds each system's responses to a common set of tasks and represents systems as points obtained by classical multidimensional scaling of the distances between their embedded responses \citep{duderstadt2023data, helm2025statistical, acharyya2024consistent}.
Most relevant to our setting, \citet{helm2026query} regress benchmark scores on DKPS representations and show that the resulting predictions are task-efficient relative to the sample score.
\trubric{} extends this geometry to agentic traces and combines it with a two-parameter IRT model, exploiting the complementary signal in traces and scores.

\section{Discussion}
\label{sec:discussion}
% TODO.

\bibliography{references}
\bibliographystyle{iclr2027_conference}

\end{document}
% build-hook smoke test
